Three extended falsification checks probe the construct validity of the AI proxy, the direction of the AI-TFP relationship, and the sensitivity of the result to the small complete-case sample.

First, replacing the composite AI index with a digital-infrastructure-only sub-index (internet use, mobile subscriptions, and secure servers) yields a coefficient of 0.075 (p-value 0.002, n = 1020). An innovation-only sub-index (resident and non-resident patent applications and IP receipts) yields 0.002 (p-value 0.728, n = 953). Both sub-indices are positively associated with TFP, indicating that the headline association is not isolated to AI-specific activity.

Second, a reverse-causality check regresses log AI adoption on one-period-lagged log TFP. The estimated coefficient is 2.306 (p-value 0.037, n = 342).

Third, restricting the baseline sample to countries with at least 8 of 15 years of non-missing AI data (n = 278 observations) yields an AI coefficient of 0.030 (p-value below 0.001).

For context, only 98 of the 193 countries in the cleaned panel report AI data in at least one year, contributing 630 country-year observations with non-missing AI data, against 2265 country-years without it. Table~\ref{tab:sample-selection} compares observable characteristics across these two groups.